% GENERATED FILE. Edit canonical lessons, then run npm run generate:books.
% canonical-source-hash: fnv1a64:72ad09c54beb3ec3
% canonical-lessons: SA-C240-tavat, SA-C240-sanaih, SA-C240-pratidinam, SA-C240-pratyaham, SA-C240-nityam, SA-R240-first-pass-words-from-chapters-232-235, SA-R240-second-pass-words-from-chapters-236-240

\chapter{So long, Gradually, Every day, Daily, Constantly}
\label{ch:sa-so-long-gradually-every-day-daily-constantly}
\hlchaptermodality{240}

\begin{chapteropening}
\textbf{By the end of this chapter:} \emph{I can say so long, gradually, every day, daily and constantly.}

Complete the last lesson of chapter 240: I can say so long, gradually, every day, daily and constantly.
\end{chapteropening}

\section[\texorpdfstring{\sk{तावत्} (tāvat)}{तावत् (tāvat)}]{\sk{तावत्} (tāvat) — so long, meanwhile}

\label{lesson:SA-C240-tavat}



\begin{quote}
Before the new one: say the Sanskrit for compassion, then the Sanskrit for pity.
\end{quote}

\subsection*{You'll want to know: \sk{तावत्}}
\textbf{\sk{तावत्}} — \emph{tāvat} — \textquotedblleft{}so long, meanwhile\textquotedblright{}.

\begin{grammarlens}[title={What you've built}]
One more word for this chapter.
\end{grammarlens}

\subsection*{Your turn}
\begin{itemize}
\raggedright
  \item \emph{Say it:} \emph{tāvat}
  \item \emph{Say it:} \emph{tāvat}, once more
  \item \emph{Say it:} say \emph{tāvat}
  \item \emph{Recall:} say the Sanskrit for compassion, then the Sanskrit for pity, then say \emph{tāvat} again
\end{itemize}

\begin{tcolorbox}[breakable,colback=teal!4,colframe=teal!35!black,title={Before you move on}]
What does \sk{तावत्} mean? (\textquotedblleft{}So long, meanwhile\textquotedblright{}.) Say it once more.
\end{tcolorbox}
\section[\texorpdfstring{\sk{शनैः} (śanaiḥ)}{शनैः (śanaiḥ)}]{\sk{शनैः} (śanaiḥ) — gradually, slowly}

\label{lesson:SA-C240-sanaih}



\begin{quote}
Before the new one: say the Sanskrit for pity, then the Sanskrit for so long.
\end{quote}

\subsection*{You'll want to know: \sk{शनैः}}
\textbf{\sk{शनैः}} — \emph{śanaiḥ} — \textquotedblleft{}gradually, slowly\textquotedblright{}.

\begin{grammarlens}[title={What you've built}]
One more word for this chapter.
\end{grammarlens}

\subsection*{Your turn}
\begin{itemize}
\raggedright
  \item \emph{Say it:} \emph{śanaiḥ}
  \item \emph{Say it:} \emph{śanaiḥ}, once more
  \item \emph{Say it:} say \emph{śanaiḥ}
  \item \emph{Recall:} say the Sanskrit for pity, then the Sanskrit for so long, then say \emph{śanaiḥ} again
\end{itemize}

\begin{tcolorbox}[breakable,colback=teal!4,colframe=teal!35!black,title={Before you move on}]
What does \sk{शनैः} mean? (\textquotedblleft{}Gradually, slowly\textquotedblright{}.) Say it once more.
\end{tcolorbox}
\section[\texorpdfstring{\sk{प्रतिदिनम्} (pratidinam)}{प्रतिदिनम् (pratidinam)}]{\sk{प्रतिदिनम्} (pratidinam) — every day}

\label{lesson:SA-C240-pratidinam}



\begin{quote}
Before the new one: say the Sanskrit for so long, then the Sanskrit for gradually.
\end{quote}

\subsection*{You'll want to know: \sk{प्रतिदिनम्}}
\textbf{\sk{प्रतिदिनम्}} — \emph{pratidinam} — \textquotedblleft{}every day\textquotedblright{}.

\begin{grammarlens}[title={What you've built}]
One more word for this chapter.
\end{grammarlens}

\subsection*{Your turn}
\begin{itemize}
\raggedright
  \item \emph{Say it:} \emph{pratidinam}
  \item \emph{Say it:} \emph{pratidinam}, once more
  \item \emph{Say it:} say \emph{pratidinam}
  \item \emph{Recall:} say the Sanskrit for so long, then the Sanskrit for gradually, then say \emph{pratidinam} again
\end{itemize}

\begin{tcolorbox}[breakable,colback=teal!4,colframe=teal!35!black,title={Before you move on}]
What does \sk{प्रतिदिनम्} mean? (\textquotedblleft{}Every day\textquotedblright{}.) Say it once more.
\end{tcolorbox}
\section[\texorpdfstring{\sk{प्रत्यहम्} (pratyaham)}{प्रत्यहम् (pratyaham)}]{\sk{प्रत्यहम्} (pratyaham) — daily}

\label{lesson:SA-C240-pratyaham}



\begin{quote}
Before the new one: say the Sanskrit for gradually, then the Sanskrit for every day.
\end{quote}

\subsection*{You'll want to know: \sk{प्रत्यहम्}}
\textbf{\sk{प्रत्यहम्}} — \emph{pratyaham} — \textquotedblleft{}daily\textquotedblright{}.

\begin{grammarlens}[title={What you've built}]
One more word for this chapter.
\end{grammarlens}

\subsection*{Your turn}
\begin{itemize}
\raggedright
  \item \emph{Say it:} \emph{pratyaham}
  \item \emph{Say it:} \emph{pratyaham}, once more
  \item \emph{Say it:} say \emph{pratyaham}
  \item \emph{Recall:} say the Sanskrit for gradually, then the Sanskrit for every day, then say \emph{pratyaham} again
\end{itemize}

\begin{tcolorbox}[breakable,colback=teal!4,colframe=teal!35!black,title={Before you move on}]
What does \sk{प्रत्यहम्} mean? (\textquotedblleft{}Daily\textquotedblright{}.) Say it once more.
\end{tcolorbox}
\section[\texorpdfstring{\sk{नित्यम्} (nityam)}{नित्यम् (nityam)}]{\sk{नित्यम्} (nityam) — constantly, always}

\label{lesson:SA-C240-nityam}



\begin{quote}
Before the new one: say the Sanskrit for every day, then the Sanskrit for daily.
\end{quote}

\subsection*{You'll want to know: \sk{नित्यम्}}
\textbf{\sk{नित्यम्}} — \emph{nityam} — \textquotedblleft{}constantly, always\textquotedblright{}.

\begin{grammarlens}[title={What you've built}]
One more word for this chapter.
\end{grammarlens}

\subsection*{Your turn}
\begin{itemize}
\raggedright
  \item \emph{Say it:} \emph{nityam}
  \item \emph{Say it:} \emph{nityam}, once more
  \item \emph{Say it:} say \emph{nityam}
  \item \emph{Recall:} say the Sanskrit for every day, then the Sanskrit for daily, then say \emph{nityam} again
\end{itemize}

\begin{tcolorbox}[breakable,colback=teal!4,colframe=teal!35!black,title={Before you move on}]
What does \sk{नित्यम्} mean? (\textquotedblleft{}Constantly, always\textquotedblright{}.) Say it once more.
\end{tcolorbox}
\section[(dialogue)]{Words from chapters 232-235 — a first pass}

\label{lesson:SA-R240-first-pass-words-from-chapters-232-235}



\begin{quote}
Say the words for daily and constantly.
\end{quote}

\subsection*{Your turn}
Say these aloud:

\begin{itemize}
\raggedright
  \item is present, is born, dies, is lost, gets angry
  \item rains, sprouts, waits, follows, goes out
  \item dries up, floats, begins, finishes, hands over
  \item covers, plants, drives, embraces, kisses
\end{itemize}

\begin{tcolorbox}[breakable,colback=teal!4,colframe=teal!35!black,title={Before you move on}]
Is present? (\textbf{\sk{वर्तते}}, \emph{vartate}.) Is born? (\textbf{\sk{जायते}}, \emph{jāyate}.) Dies? (\textbf{\sk{म्रियते}}, \emph{mriyate}.) Is lost? (\textbf{\sk{नश्यति}}, \emph{naśyati}.) Gets angry? (\textbf{\sk{कुप्यति}}, \emph{kupyati}.) Rains? (\textbf{\sk{वर्षति}}, \emph{varṣati}.) Sprouts? (\textbf{\sk{रोहति}}, \emph{rohati}.) Waits? (\textbf{\sk{प्रतीक्षते}}, \emph{pratīkṣate}.) Follows? (\textbf{\sk{अनुगच्छति}}, \emph{anugacchati}.) Goes out? (\textbf{\sk{निर्गच्छति}}, \emph{nirgacchati}.) Dries up? (\textbf{\sk{शुष्यति}}, \emph{śuṣyati}.) Floats? (\textbf{\sk{प्लवते}}, \emph{plavate}.) Begins? (\textbf{\sk{आरभते}}, \emph{ārabhate}.) Finishes? (\textbf{\sk{समापयति}}, \emph{samāpayati}.) Hands over? (\textbf{\sk{प्रयच्छति}}, \emph{prayacchati}.) Covers? (\textbf{\sk{पिदधाति}}, \emph{pidadhāti}.) Plants? (\textbf{\sk{आरोपयति}}, \emph{āropayati}.) Drives? (\textbf{\sk{चालयति}}, \emph{cālayati}.) Embraces? (\textbf{\sk{आलिंगति}}, \emph{āliṅgati}.) Kisses? (\textbf{\sk{चुम्बति}}, \emph{cumbati}.)
\end{tcolorbox}
\section[(dialogue)]{Words from chapters 236-240 — a second pass}

\label{lesson:SA-R240-second-pass-words-from-chapters-236-240}



\begin{quote}
Say the words for a baby and a boy.
\end{quote}

\subsection*{Your turn}
Say these aloud:

\begin{itemize}
\raggedright
  \item a baby, a boy, a maiden, a young woman, a person
  \item a potter, a goldsmith, a blacksmith, a weaver, a washerman
  \item a nurse, an engineer, a scientist, a policeman, a driver
  \item hope, love, affection, compassion, pity
  \item so long, gradually, every day, daily, constantly
\end{itemize}

\begin{tcolorbox}[breakable,colback=teal!4,colframe=teal!35!black,title={Before you move on}]
A baby? (\textbf{\sk{शिशुः}}, \emph{śiśuḥ}.) A boy? (\textbf{\sk{कुमारः}}, \emph{kumāraḥ}.) A maiden? (\textbf{\sk{कुमारी}}, \emph{kumārī}.) A young woman? (\textbf{\sk{युवतिः}}, \emph{yuvatiḥ}.) A person? (\textbf{\sk{पुरुषः}}, \emph{puruṣaḥ}.) A potter? (\textbf{\sk{कुम्भकारः}}, \emph{kumbhakāraḥ}.) A goldsmith? (\textbf{\sk{स्वर्णकारः}}, \emph{svarṇakāraḥ}.) A blacksmith? (\textbf{\sk{लोहकारः}}, \emph{lohakāraḥ}.) A weaver? (\textbf{\sk{तन्तुवायः}}, \emph{tantuvāyaḥ}.) A washerman? (\textbf{\sk{रजकः}}, \emph{rajakaḥ}.) A nurse? (\textbf{\sk{परिचारिका}}, \emph{paricārikā}.) An engineer? (\textbf{\sk{अभियन्ता}}, \emph{abhiyantā}.) A scientist? (\textbf{\sk{वैज्ञानिकः}}, \emph{vaijñānikaḥ}.) A policeman? (\textbf{\sk{आरक्षकः}}, \emph{ārakṣakaḥ}.) A driver? (\textbf{\sk{चालकः}}, \emph{cālakaḥ}.) Hope? (\textbf{\sk{आशा}}, \emph{āśā}.) Love? (\textbf{\sk{प्रेम}}, \emph{prema}.) Affection? (\textbf{\sk{स्नेहः}}, \emph{snehaḥ}.) Compassion? (\textbf{\sk{दया}}, \emph{dayā}.) Pity? (\textbf{\sk{करुणा}}, \emph{karuṇā}.) So long? (\textbf{\sk{तावत्}}, \emph{tāvat}.) Gradually? (\textbf{\sk{शनैः}}, \emph{śanaiḥ}.) Every day? (\textbf{\sk{प्रतिदिनम्}}, \emph{pratidinam}.) Daily? (\textbf{\sk{प्रत्यहम्}}, \emph{pratyaham}.) Constantly? (\textbf{\sk{नित्यम्}}, \emph{nityam}.)
\end{tcolorbox}
